LiDAR-Inertial-Visual (LIV) fusion has been extensively validated in
classical SLAM~\cite{xu2022fastlio2,zheng2024_fastlivo2,lin2024_r3live},
providing robust metric-scale odometry through tightly-coupled iterated
extended Kalman filters (IESKF) or continuous-time factor-graph optimization.
Integrating 3DGS into this framework creates systems that combine precise
pose estimation with photo-realistic dense mapping.  We categorize the
literature into \emph{tightly-coupled} and \emph{decoupled} approaches
based on whether the Gaussian map feeds back into the odometry.

\subsection{Gaussian-LIC and Gaussian-LIC2}

\textbf{Gaussian-LIC}~\cite{lang2025_gaussianlic} (ICRA 2025) is the first
real-time photo-realistic SLAM system combining 3DGS with
LiDAR-Inertial-Camera fusion.  It uses a continuous-time trajectory
optimization framework for robust pose estimation, initializes Gaussians
from both LiDAR points and triangulated visual SfM points, and models sky
and varying camera exposure for outdoor rendering.  The system is
implemented entirely in C++/CUDA with carefully designed acceleration
strategies.  \textbf{Code}: \url{https://github.com/APRIL-ZJU/Gaussian-LIC}.

\textbf{Gaussian-LIC2}~\cite{lang2026_gaussianlic2} (IJRR 2026) extends this
to a journal version addressing three key issues: (1)~under-reconstruction
in LiDAR-blind areas, solved via a lightweight zero-shot depth completion
network that fuses RGB appearance cues with sparse LiDAR measurements;
(2)~geometric accuracy, enhanced by using sparse but precise LiDAR depths to
supervise Gaussian optimization with CUDA-accelerated strategies;
(3)~bidirectional coupling between odometry and mapping, where photometric
constraints from the incrementally built Gaussian map are fed back into
continuous-time factor graph optimization to improve tracking under LiDAR
degeneration.  A dedicated dataset with ground-truth poses, depths, and
extrapolated trajectories is released for evaluating out-of-sequence novel
view synthesis.  \textbf{Code}: \url{https://xingxingzuo.github.io/gaussian_lic2}.

\subsection{GS-LIVM}

GS-LIVM~\cite{xie2024_gslivm} (ICCV 2025 Workshop) targets real-time
photo-realistic LiDAR-Inertial-Visual mapping for large-scale outdoor
scenes.  Its key innovation is the use of Gaussian Process Regression (GPR)
at the voxel level (Voxel-GPR) to densify sparse and unevenly distributed
LiDAR observations into uniform point grids.  The estimated covariance from
GPR initializes the scale and rotation of 3D Gaussians, accelerating
convergence.  Custom CUDA kernels enable real-time performance.  Evaluated
on outdoor datasets including the Botanic Garden sequence ($\sim$200\,m
trajectory).  \textbf{Code}: \url{https://github.com/xieyusen/GS-LIVM}.

\subsection{GS-LIVO}

GS-LIVO~\cite{hong2025_gslivo} (IEEE T-RO 2025) proposes a global Gaussian
map with hash-indexed voxels organized in a recursive octree, adapting to
different environment scales.  It uses an IESKF for tightly-coupled
multi-sensor fusion and maintains a sliding window of Gaussians to reduce
GPU memory.  Notably, it is the first Gaussian-based SLAM deployable on
embedded systems (NVIDIA Jetson Orin NX), demonstrating real-time performance
on resource-constrained platforms---a critical milestone for robotic
deployment.

\subsection{LVI-GS}

LVI-GS~\cite{zhao2025_lvigs} (IEEE T-IM 2025) introduces a
tightly-coupled LiDAR-Visual-Inertial mapping framework.  Gaussians are
initialized from colorized LiDAR points and optimized via differentiable
rendering.  A pyramid-based training approach learns multi-level features,
and a depth loss derived from LiDAR measurements improves geometric
perception.  Thread management and custom CUDA acceleration enable real-time
operation.  \textbf{Project page}: \url{https://kwanwaipang.github.io/LVI-GS/}.

\subsection{RMGS-SLAM}

RMGS-SLAM~\cite{li2026_rmgs} (arXiv 2026) is a tightly-coupled LIV 3DGS-SLAM
with explicit loop closure.  Its cascaded Gaussian initialization combines
feed-forward predictions with voxel-PCA geometric priors, improving
initialization quality and convergence speed.  Loop closure is performed
directly on the optimized Gaussian map via Gaussian-based GICP registration,
followed by pose-graph optimization.  A new benchmark with
hardware-synchronized solid-state LiDAR, camera, and IMU plus ground truth
is released.  The system achieves a strong balance among real-time
efficiency, localization accuracy, and rendering quality.

\subsection{Real-Time LiDAR Gaussian Splatting SLAM}

Tak et al.~\cite{tak2026_realtimelidar} (arXiv 2026) couple fast G-ICP
registration with spherical rasterization-based dense mapping.  A
covariance-derived geometry score measures local complexity and drives
pruning in planar regions and selective densification in structurally rich
areas, enabling real-time large-scale mapping.

\subsection{LV-GS SLAM}

LV-GS SLAM~\cite{he2026lvgsslam} (Applied Sciences 2026) integrates LiDAR
and visual data for large-scale incremental reconstruction with real-time
tracking.  A LiDAR odometry front-end processes 30\,Hz LiDAR at 15--19\,Hz
tracking.  A novel depth propagation approach initializes 3D Gaussians from
dense depth maps, achieving 9$\times$ faster PSNR convergence versus direct
LiDAR initialization.  Keyframe-based submap management dynamically adjusts
memory allocation to prevent GPU overflow.  Validated on KITTI and a
quadruped robot platform.

\subsection{LIT-GS: LiDAR-Inertial-Thermal}

LIT-GS~\cite{shi2026_litgs} (IROS 2026) extends the LIV paradigm to thermal
imaging, replacing the RGB camera with a long-wave thermal sensor for
illumination-robust mapping.  It injects LiDAR-derived plane geometry as an
explicit constraint in both pose/structure refinement (via weighted
point-to-plane BA) and Gaussian optimization (via a LiDAR-plane-regularized
splatting objective).  This addresses the weak texture and low contrast of
thermal images, which would otherwise corrupt photometric supervision.

\subsection{Summary of LiDAR-Based Systems}

\Cref{tab:lidar_summary} compares the key characteristics of LiDAR+RGB+IMU
Gaussian SLAM systems.  The evolution shows clear trends: (1)~from
offline to real-time operation, (2)~from LiDAR-only initialization to
hybrid LiDAR+depth-completion initialization, (3)~from decoupled
odometry-mapping to bidirectional coupling where the Gaussian map improves
tracking, and (4)~from desktop GPUs to embedded deployment.

\begin{table}[t]
\centering
\caption{LiDAR + RGB + IMU Gaussian Splatting SLAM systems.}
\label{tab:lidar_summary}
\scriptsize
\begin{tabular}{lccccc}
\toprule
\textbf{System} & \textbf{Year} & \textbf{Venue} & \textbf{Real-time} & \textbf{Loop Closure} & \textbf{Code} \\
\midrule
Gaussian-LIC~\cite{lang2025_gaussianlic} & 2025 & ICRA & \checkmark & -- & \checkmark \\
GS-LIVM~\cite{xie2024_gslivm} & 2024 & arXiv & \checkmark & -- & \checkmark \\
GS-LIVO~\cite{hong2025_gslivo} & 2025 & T-RO & \checkmark & -- & -- \\
LVI-GS~\cite{zhao2025_lvigs} & 2025 & T-IM & \checkmark & -- & -- \\
Gaussian-LIC2~\cite{lang2026_gaussianlic2} & 2026 & IJRR & \checkmark & -- & \checkmark \\
RMGS-SLAM~\cite{li2026_rmgs} & 2026 & arXiv & \checkmark & \checkmark & -- \\
RT-LiDAR-GS~\cite{tak2026_realtimelidar} & 2026 & arXiv & \checkmark & -- & -- \\
LV-GS SLAM~\cite{he2026lvgsslam} & 2026 & AppSci & \checkmark & -- & -- \\
LIT-GS~\cite{shi2026_litgs} & 2026 & IROS & \checkmark & -- & -- \\
\bottomrule
\end{tabular}
\end{table}

\textbf{Key insight}: The dominant challenge in LiDAR-based Gaussian SLAM is
\emph{LiDAR sparsity and FoV mismatch with the camera}.  Gaussian-LIC2's
depth completion network and RMGS-SLAM's cascaded initialization are the two
most effective solutions, while GS-LIVM's Voxel-GPR densification provides a
lightweight alternative.  The bidirectional coupling of Gaussian-LIC2
(map$\rightarrow$odometry photometric feedback) is a unique contribution
that demonstrates the Gaussian map is not merely a passive visualization
tool but an active constraint source for tracking.
